\documentclass[11pt]{article}
\usepackage[margin=1in]{geometry}
\usepackage{amsmath,amssymb,amsthm}
\usepackage[hidelinks]{hyperref}
\newtheorem{theorem}{Theorem}
\title{An effective quadratic lower bound for unweighted Dominating Set}
\author{Atlas}
\date{18 November 2026}
\begin{document}
\maketitle

A dominating set of a graph $G$ is a set $D\subseteq V(G)$ such that every
vertex outside $D$ has a neighbor in $D$. The input to \textsc{Dominating
Set} is an explicit simple undirected graph and a nonnegative binary budget.
The question is whether there is a dominating set of size at most the
budget. A graph is monopolar if its vertex set is the union of an
independent set and a disjoint union of cliques.

For $X\subseteq V(G)$, the cut rank is the rank over $\mathbb F_2$ of
the adjacency matrix between $X$ and $V(G)\setminus X$. The width of a
vertex order is the largest cut rank of any prefix, and linear rank-width
is the minimum such width. A rank decomposition is a subcubic tree whose
leaves correspond bijectively to the vertices. Each edge induces a cut,
and the width of the decomposition is the largest cut rank of those cuts.
Rank-width is the minimum width of such a tree. A supplied order or
decomposition is measured by its actual width.

For a computable nonnegative integer function $h$, write
$h(p)=o_{\mathrm{eff}}(p^2)$ when there is a computable function $T$
such that $h(p)\le\eta\max(1,p)^2$ for all integers $p\ge T(\eta)$,
for each positive rational $\eta$.

Liu and Meng~\cite[Theorem~3.1 and Corollary~3.10]{lm} prove a quadratic
lower bound for unweighted \textsc{Dominating Set} under ETH. The theorem
below is a direct consequence of their monopolar construction in
Section~3.2. Appendix~C of~\cite{lm} observes that the construction
works for clauses of unrestricted length. This lets us start from a
CNF formula encoding a Multicolored Clique witness with $O(\log n)$ bits
per color, where $n$ is the source graph order. The resulting width is
linear in the square root of the number of formula variables, including
for the supplied order and decomposition. The effective convergence
cutoff then turns a hypothetical solver into a fixed-parameter algorithm
for Multicolored Clique.

\hypertarget{dominating-set-rankwidth}{}
\begin{theorem}\label{dominating-set-rankwidth}
Assume deterministic $\mathrm{FPT}\ne\mathrm W[1]$. For every fixed
$C\ge0$ and every computable $h(p)=o_{\mathrm{eff}}(p^2)$, no deterministic
algorithm decides \textsc{Dominating Set} in $O(2^{h(p)}N^C)$ time on
monopolar graphs. This holds for rank-width and linear rank-width, and
separately for the width of a supplied rank decomposition or vertex order.
Here $N$ is the complete binary input length, including the certificate
when supplied.
\end{theorem}

\section{A compact formula for Multicolored Clique}

Let a Multicolored Clique instance have $\ell$ colors and $n\ge2$ vertices.
We may decide cases with empty colors or fewer than two colors directly.
Use a block of $b=\lceil\log_2 n\rceil$ variables to encode one vertex of
each color. Forbid unused codes by adding one clause for each unused code.
For every nonedge with endpoints of different colors add a clause forbidding
the pair of codes. A forbidden bit pattern is represented by the disjunction
of the literals that disagree with it. A satisfying assignment is therefore
equivalent to a multicolored clique. This uses $\ell b$ variables and at most
$\ell 2^b+\ell^2 n^2$ clauses of length at most $2b$. Both the construction
and the full encoding take $F(\ell)n^{O(1)}$ time for a computable $F$.

An empty conjunction is satisfiable and can be decided directly. Let
$C_1,\ldots,C_m$ denote the clauses, where $m\ge1$. Pad with unused
variables to obtain a $k\times k$ array of variables with
$k=\lceil\sqrt{\ell b}\rceil$. An assignment is now a binary matrix
whose $a$th row assigns the variables in row $a$.

\section{The reduction to domination}

Use one layer per clause. For each layer $i$ and row $a$ create a clique
$A_{i,a}$ with one vertex $v_{i,a,x}$ for each $x\in\mathbb F_2^k$.
Add two nonadjacent vertices adjacent to every vertex of $A_{i,a}$ and
with no other neighbors. These force a choice from each clique at the
budget $km$. In layer $i$ add a clause vertex adjacent to $v_{i,a,x}$
exactly when $x$ satisfies a literal of $C_i$ from row $a$.

To make consecutive layer assignments agree, for each $i<m$ and
each $t,p,r\in\mathbb F_2^k$ with $r\ne0$ add a vertex $e_{i,t,p,r}$.
Join it to $v_{i,a,x}$ if $x\cdot t\ne p_a$, and to $v_{i+1,a,x}$ if
$x\cdot t\ne p_a+r_a$. There are no further edges. The choice cliques
and the independent set of remaining vertices give a monopolar partition.

Every dominating set contains at least one vertex per choice clique and
its two guards taken together. If it contains no choice vertex in such
a group, it must contain both guards. Thus a solution of size at most
$km$ selects exactly one vertex from every choice clique and no other
vertices. Write $X$ and $Y$ for the selected assignments at two consecutive
layers. The test vertex indexed by $(t,p,r)$ is undominated exactly when
$Xt=p$ and $Yt=p+r$. If $X=Y$, this never occurs. If $X\ne Y$, choose a
unit vector $t$ at a column on which they differ, and let $p=Xt$ and
$r=(X+Y)t$. The resulting test vertex is undominated. Hence all layer
assignments agree in a solution. Every clause vertex must have a selected
neighbor, so this common assignment satisfies the formula. Conversely, a
satisfying assignment used in all layers dominates the choice cliques,
guards, clause vertices and test vertices.

\section{Width and running time}

Order the vertices in layer order, with the tests between layers $i$ and
$i+1$ in one block between the layers. Within each layer put its clause
vertex first, followed by consecutive row groups with the guards before
the assignment vertices. Across a left test interface the entry for
$(a,x)$ and $(t,p,r)$ is $x\cdot t+p_a$. Every row is in the span of the
$2k$ coordinate functions $t_j$ and $p_a$. Across a right interface use
$t_j$ and $p_a+r_a$ in the same way. The rank of either interface, or any
submatrix, is therefore at most $2k$.

A cut inside a test block splits only its two test interfaces. A cut
inside a layer splits at most the preceding and succeeding test interfaces,
one row group and the clause adjacency. For the split row group, every
nonzero row toward the unprocessed portion of the clique is the all-ones
row, and the unprocessed guards have no processed clique vertex preceding
them. This adds at most one to the rank, and the clause vertex adds at most
one. Thus the supplied order has width at most $4k+2$. A rank decomposition
of width at most $4k+2$ follows by placing one leaf per vertex along a
path, since its remaining cuts isolate a vertex. Both intrinsic widths
also satisfy this bound.

The layers use $mk(2^k+2)+m$ vertices and the intervening blocks use
$(m-1)2^{2k}(2^k-1)$. The complete explicit graph, budget and either
certificate can be constructed in $2^{O(k)}$ times a fixed polynomial in
the formula length. This bound allows dense edge output. Because
$\sqrt{\ell\log_2 n}\le(\ell+\log_2 n)/2$, the construction takes
$F'(\ell)n^{O(1)}$ time for a computable $F'$. The output parameter
satisfies $\max(1,p)^2\le A(\ell)\log_2 n$ for a computable positive
integer $A(\ell)$.

Suppose the algorithm from the theorem exists. At
$p\ge T(1/A(\ell))$ we have $2^{h(p)}\le n$. Below the cutoff its
contribution has a computable parameter-only bound obtained by taking the
maximum over the finite prefix of $h$. The complete composition then
solves Multicolored Clique in deterministic fixed-parameter time by
number of colors. This contradicts its $\mathrm W[1]$-hardness, which
follows by copying each vertex into each color and joining distinct-color
copies exactly when their originals are adjacent~\cite{df}.

\paragraph{Acknowledgment.}
The proof was found using ChatGPT.

\begin{thebibliography}{9}
\bibitem{df}
R. G. Downey and M. R. Fellows.
Fixed-parameter tractability and completeness II: On completeness for
$\mathrm W[1]$. \emph{Theoretical Computer Science}, 141(1--2):109--131,
1995. \href{https://doi.org/10.1016/0304-3975(94)00097-3}{DOI}.
\bibitem{lm}
C. Liu and B. Meng.
Lower bounds for domination-type problems parameterized by rank-width.
\href{https://arxiv.org/abs/2608.18854v1}{arXiv:2608.18854v1}, 2026.
See Section 3.2 and Appendix C for the construction and its arbitrary-CNF extension.
\end{thebibliography}
\end{document}
